% !PRINTED 627
\weberpage{627}
b) Ist $p-1$ nicht durch 4 theilbar, also
\weq{11}{p \equiv -5^\beta \, (\mathrm{mod}\ m),}
so ist $f$ die kleinste positive Zahl, die den Congruenzen
\[ f \equiv 0 \, (\mathrm{mod}\ 2), \quad 2\beta f \equiv 0 \, [\mathrm{mod}\ \varphi(m)] \]
genügt, und folglich ist $e$ der grösste gemeinschaftliche Theiler von $\varphi(m)$ und $2\beta$, oder, wenn $2\beta$ durch $\varphi(m)$ theilbar ist, gleich $\tfrac{1}{2}\varphi(m)$. (Eine Ausnahme bildet hier der Fall $m=4$, in dem $e=1$ ist. Diesen Fall, der ja schon in früheren Abschnitten vollständig erledigt ist, lassen wir daher jetzt bei Seite.)

Es ergiebt sich dann aus (11), dass eine Zahl von der Form $5^\lambda$ nur dann einer Potenz von $p$ congruent sein kann, wenn $\lambda$ ein Vielfaches von $2\beta$, also ein Vielfaches von $e$ ist.

Hier kann nun in jeder der Nebengruppen $\mathfrak{A}\xi$ eine Zahl gefunden werden, die $\equiv 1 \, (\mathrm{mod}\ 4)$ ist, da $p \equiv -1 \, (\mathrm{mod}\ 4)$ eine Zahl in $\mathfrak{A}$ ist, und wir erhalten demnach ein vollständiges System der $\xi$ in der Form:
\[ 1, 5, 5^2, \ldots, 5^{e-1}, \]
und darin ist $e$ höchstens $= \tfrac{1}{2}\varphi(m) = 2^{\varkappa-2}$.

Wir fassen das Bewiesene im folgenden Theorem zusammen:

III. \gesperrt{Zerlegt man die von $q$ verschiedene Primzahl $p$ im Körper $\Omega_m$ in ihre Primfactoren:}
\[ p = \mathfrak{p}_{\xi_1} \mathfrak{p}_{\xi_2} \cdots \mathfrak{p}_{\xi_e}, \]
\gesperrt{so kann man, wenn $p-1$ durch $q$, oder im Falle eines geraden $m$ durch 4 theilbar ist, die $\xi$ ein volles System durch $q$ nicht theilbarer Reste nach dem Modul $m_1$ durchlaufen lassen, und die $\mathfrak{p}_\xi$ als Primfunctionale im Körper $\Omega_{m_1}$ annehmen, wenn $m_1$ der grösste gemeinschaftliche Theiler von $p-1$ und $m$ ist.}

\gesperrt{Ist aber $p-1$ nicht durch $q$ oder für ein gerades $m$ nicht durch 4 theilbar, und bedeutet $g$ im ersten Falle eine Primitivwurzel von $m$, im zweiten Falle die Zahl 5, so durchläuft $\xi$ die Reihe der Zahlen}
\[ 1, g, g^2, \ldots, g^{e-1}, \]
\gesperrt{und $e$ ist bei einem ungeraden $m$ nicht durch $q-1$ theilbar, und bei geradem $m$ mindestens $=2$ und höchstens $=\tfrac{1}{2}\varphi(m)$. Ausgeschlossen ist hierbei der Fall $m=4$.}
